\section{Related Work}
\label{sec:related}

\paragraph{Spectrum-based fault localisation.}
SBFL methods rank statements or functions by suspiciousness scores derived from
pass/fail test-coverage data.
Tarantula~\cite{jones05tarantula} introduced the coefficient refined by many
later techniques.
A broad survey is given by \citet{wong16survey}.
Our work differs fundamentally: we use neither test suites nor execution coverage;
the LLM reasons directly from a natural-language symptom description and optional
call-graph context.

\paragraph{LLMs for code evaluation and program repair.}
LLMs have been extensively evaluated on code generation and automated program
repair~\cite{sobania23chatgpt,liu23evalplus}.
While benchmark evaluations such as EvalPlus~\cite{liu23evalplus} assess standalone
generation correctness, automated repair tools prompt models iteratively with compiler
feedback or test outcomes.
Our study focuses on the upstream task of identifying the faulty function
prior to repair, investigating whether call-graph representations provide
sufficient topological guidance to isolate root causes in complex C software.

\paragraph{Repository navigation agents and context engineering.}
Recent work investigates how large language models can navigate repository-scale
codebases using structured tools.
Most closely related to our agent architecture, \citet{yildiz25jitvul} introduce
the JitVul benchmark, deploying ReAct-style agents equipped with interprocedural
context to detect vulnerabilities in C repositories.
In an empirical investigation of context engineering, \citet{li26onesize}
demonstrate that ``one size does not fit all'' when retrieving repository context:
while navigation-based agents dynamically explore dependencies, navigation-based
approaches incur 10--20$\times$ higher cost.
Our empirical findings directly mirror this cost trade-off, showing a 10--23$\times$
prompt-token multiplier for tool-navigating agents over static BFS text slices, while
providing an edge-level audit of the structural gaps and cited edge validity
inherent to static call graphs on large C systems.

\paragraph{LLM agents with structured tool use.}
The \textsc{ReAct} framework~\cite{yao23react} established the
reason-act-observe pattern for autonomous tool-using agents.
We apply this pattern to static call-graph exploration, ablating the available
tool primitives (call chains, type hierarchies, taint paths, raw source retrieval)
to isolate which specific tools assist localisation and whether tool access reduces
the citing of edges not supported by the underlying call graph.
